\section{Are You a Counterrevolutionary?}

\label{sec:Counterrevolutionary}

\begin{quotex}
What is needed is not a revolution in the opposite direction, but the opposite of a revolution.
\flright{\textsc{Joseph de Maistre}}
\end{quotex}


\begin{quotex}
My principles are only those that prior to the French revolution, every well born person considered healthy and normal.
\flright{\textsc{Julius Evola}}
\end{quotex}


At a time when the right is in a state of confusion, we witness many attempts at its definition. The term ``conservative'', with all its various qualifications, doesn't satisfy. But ``right'' is vague, meaning little more than ``not left'', despite attempts to define a ``new right'', or an ``alternative right'', as though there is a mainstream right.

No, the correct term is ``counterrevolution'', that is, the restoration of order that every succeeding revolution has sought to undermine. As Evola points out, at one time every well born man was ``on the right'', so the distinction was hardly necessary until the aftermath of the French revolution.

I have prepared a little quiz for readers to determine if they are on the side of the revolution or on the counterrevolution. For each position in the table below, circle either the option on the left, or the one on the right. Score 5 points for every circle on the right column. Score 0 points if you circle any option on the left. Your score indicates how healthy and normal is your point of view. If you circled any on the left, the count will indicate whether you are in the revolutionary vanguard or merely a useful idiot.

\begin{small}
\begin{longtable}{@{}rp{.25\textwidth}p{.3\textwidth}@{}}
\caption{Counterrevolutionary Evaluation Form}\\

\toprule

~ & \textit{Left} & \textit{Right} \\
\midrule

\endfirsthead

\toprule

~ & LEFT & RIGHT \\
\midrule

\endhead

\bottomrule

\endlastfoot

Foundation of Reality & Matter & Spirit \\ Philosophy & Nominalism & Realism (ideas) \\ Causation is & Horizontal & Horizontal and Vertical \\ Reality is & Socially constructed & Reflection of Logos \\ Change is the result of & Evolution & Creativity \\ Social change is caused by & Historical or random forces & Conscious agents \\ Basic societal unit & Individual & Family \\ Societal structure & Egalitarian & Hierarchical \\ Group interaction & Class/ethnic/gender warfare & Cooperative class structure \\ Relations are & Contractual, consensual & Organic \\ Highest authority & State & Transcendence \\ Function of the state & Welfare State & Ethical State \\ Evil derives from & Political or social institutions & Man himself \\ Knowledge is & Opinion (knowledge of the particular) & Wisdom (knowledge of the idea) \\ Social Justice & To each according to his need & To each his just desserts \\ Equality of & The Individual or Results & The person \\ Man is a & Body & Soul \\ Man's mind consists of & Blank slate & Inherited characteristics \\ Freedom to act on & Instinctual urges & Rational desires \\
\end{longtable}\end{small}



\flrightit{Posted on 2010-03-09 by Cologero}

\begin{center}* * *\end{center}

\begin{footnotesize}
\begin{sffamily}
\texttt{GFJF on 2010-08-25 at 11:44 said:}

I wonder if it isn't unhelpful to call these positions left and right. Mightn't `technical' and `organic' be more accurate?

\hfill

\texttt{GFJF on 2010-08-25 at 13:08 said:}

To anticipate, I don't think this is a question of mere words. The terms left and right create confusion because they are taken to mean certain things already. Thus even in our own private theorizing we have constantly to qualify different stances as `not really right' and `not really left', which causes more than necessary mental convolution. In conversation with others it fosters the impression that Tradition is simply a right-wing political movement, one of the many that arbitrarily defines itself as `the True Right'. We should keep left and right as descriptors of different stances taken *within* the mechanical worldview, and not try to alter their definitions.

\hfill

\texttt{Cologero on 2010-08-25 at 18:39 said:}

There is only confusion because people use words improperly, whether out of ignorance or to further their own ends. There is absolutely no ``arbitrariness'' to the word ``right'' in a political sense. It has a definite historical origin and meaning, referring to those defending ``throne and altar'' against the Jacobins. Tradition is not a so-called ``right-wing'' political movement; it simply opposes any non-Traditional movement or program.

Now ask yourself who indeed altered the definitions and for what ends.

\hfill

\texttt{GFJF on 2010-08-25 at 20:57 said:}

Nevertheless it *seems* like an arbitrary redefinition to most people and causes them to lump Tradition in with all the other particular movements. Even if we aren't interested in explaining ourselves to most people, we should be able to do so, and eloquently, if called on. So at the very least, technical and organic are handy elaborations of what we mean by left and right.

\hfill

\texttt{Cologero on 2010-08-25 at 22:01 said:}

I believe the explanation I gave you of ``right'' is easy to understand. Try to explain ``technical'' and ``organic'' in two sentences.

\hfill

\texttt{GFJF on 2010-08-25 at 22:18 said:}

You're right. Throne and altar says it all.

\hfill

\texttt{Mark Citadel on 2016-04-10 at 08:40 said:}

This is excellent, and could not have come down at a better time! Confusions in terms which stem from left wing misuse have led to much dissolution and damaging masquerades. Every ideology that accepts even a fraction of the essential post-Enlightenment revolutions (ALL of them), is to greater or lesser degree a viral extension of the problem. Puts what I had written in my essay on what `right' means into a very easy format.

One query, when you speak of social change being caused either by ``historical or random forces'' (left), or by ``conscious agents'' (right), I understand what you are getting at, but does this account for the cyclical change in universal patterns which tends towards either degeneration or rejuvination? Perhaps we should interpret these in light of the Vedic tradition, in the same way that we would interpret God's foreknowledge not impacting the free will of human beings, that prophesy and fate are not at odds with our free will, but rather just predict them? Is this accurate?

\hfill

\texttt{Cologero on 2016-04-11 at 20:23 said:}

RE: cycles

At any particular moment, there are possibilities of manifestation. As disorder (or entropy) increases, there are more such possibilities. Viz., things that had not been possible, are now possible. For any system, we need to look at several elements:

\begin{enumerate}


\item The supernatural source of an idea: where in the cosmological hierarchy did it originate. Or is it an idea of the anti-Tradition. These are all possibilities of manifestation.

\item The public agents of change: they promote the idea, good or bad, and do so openly.

\item The hidden agents of change: these are organizations that promote the idea without making themselves known.

\item The unconscious agents are basically everyone else who simply absorb the ideas without understanding them, their true source, or ultimate consequences.
\end{enumerate}


So it is easy enough to predict the results: Ideas have consequences and all possibilities of manifestation must indeed manifest.

\hfill

\end{sffamily}
\end{footnotesize}

